% Sección 4: Bibliografía
% Generado por latex-writer | Fecha: 2026-03-25
% =============================================================================

\section{Bibliografía}
\label{sec:bibliografia}

A continuación se relacionan las fuentes consultadas para el desarrollo del diseño termohidráulico de la instalación de vapor y condensados. Las referencias se han organizado en cinco categorías según su naturaleza: normativas y reglamentos de obligado cumplimiento, guías técnicas institucionales, manuales de eficiencia energética, catálogos comerciales de fabricantes y tablas de propiedades termodinámicas. La numeración correlativa permite su citación inequívoca a lo largo de la memoria técnica.

\subsection*{Normativas y Reglamentos}

\begin{enumerate}[label={[\arabic*]}, leftmargin=2em]
	\item \textbf{RITE --- Reglamento de Instalaciones Térmicas en los Edificios.}
	      Real Decreto 1027/2007, de 20 de julio.
	      Ministerio de Industria, Turismo y Comercio (España), 2007.
	      \textit{Criterio de temperatura superficial máxima en zonas accesibles ($T_{\text{sup}} < 30$~°C).}

	\item \textbf{ISO 12241:2022 --- Thermal insulation for building equipment and industrial installations: Calculation rules.}
	      International Organization for Standardization (ISO), 2022.
	      \textit{Metodología de cálculo de espesores de aislamiento térmico y flujos de calor.}

	\item \textbf{DIN 2448 --- Seamless steel tubes: Technical delivery conditions.}
	      Deutsches Institut für Normung (DIN).
	      \textit{Especificación de tuberías de acero sin soldadura para diámetros nominales DN~80--150.}

	\item \textbf{API Schedule (ANSI B36.10) --- Welded and Seamless Wrought Steel Pipe.}
	      American Petroleum Institute (API) / American National Standards Institute (ANSI).
	      \textit{Especificación de espesores de pared para tuberías Schedule~40, 80 y 160.}
\end{enumerate}

\subsection*{Guías Técnicas Institucionales}

\begin{enumerate}[label={[\arabic*]}, leftmargin=2em, resume]
	\item \textbf{Guía Técnica de Diseño y Cálculo del Aislamiento Térmico de Conducciones, Aparatos y Equipos.}
	      Instituto para la Diversificación y Ahorro de la Energía (IDAE).
	      \textit{Materiales aislantes, cálculo de espesores óptimos y criterios de estandarización por diámetro nominal.}

	\item \textbf{VAPOR SPIRAX SARCO --- Guía de Referencia Técnica.}
	      Spirax Sarco Limited.
	      \textit{Dimensionado de tuberías de vapor, líneas de distribución y purga, tablas de factores de presión y prevención de golpe de ariete.}

	\item \textbf{Catálogo General de Productos de Tubería Industrial.}
	      Tubacero (fabricante).
	      \textit{Dimensiones nominales, materiales, espesores y conexiones para tuberías de acero.}
\end{enumerate}

\subsection*{Manuales Técnicos}

\begin{enumerate}[label={[\arabic*]}, leftmargin=2em, resume]
	\item \textbf{Manual de Eficiencia Energética en Redes de Vapor y Condensados.}
	      Ente Regional de la Energía de Castilla y León (EREN), 2010.
	      \textit{Cálculo hidráulico mediante Darcy-Weisbach y Colebrook-White; longitudes equivalentes de accesorios ($L_e/D = 26$ para codos de radio mediano); recuperación de condensados y vapor flash; dimensionamiento de pozos de goteo y estaciones de purga.}

	\item \textbf{Vapor Interesante --- Manual de Referencia Técnica.}
	      \textit{Información práctica y teórica sobre sistemas de vapor industrial.}
\end{enumerate}

\subsection*{Catálogos Comerciales}

\begin{enumerate}[label={[\arabic*]}, leftmargin=2em, resume]
	\item \textbf{Caldera de vapor de alta presión Viessmann Vitomax 100-HS, Modelo M33A --- Datos Técnicos.}
	      Viessmann Werke GmbH \& Co. KG.
	      \textit{Especificaciones técnicas del equipo seleccionado: caudal 5400~kg/h, presión 10~bar(g), temperatura 220~°C.}

	\item \textbf{Caldera Viessmann Vitomax 200-HS, Modelo M237 --- Datos Técnicos.}
	      Viessmann Werke GmbH \& Co. KG.
	      \textit{Documentación de referencia complementaria (modelo alternativo de mayor capacidad).}
\end{enumerate}

\subsection*{Tablas Termodinámicas}

\begin{enumerate}[label={[\arabic*]}, leftmargin=2em, resume]
	\item \textbf{IAPWS-IF97 --- International Association for the Properties of Water and Steam: Industrial Formulation 1997 for the Thermodynamic Properties of Water and Steam.}
	      International Association for the Properties of Water and Steam (IAPWS), 1997 (con revisiones posteriores).
	      \textit{Propiedades termofísicas del vapor y agua: entalpías, densidades, viscosidades, temperaturas de saturación, volúmenes específicos.}
\end{enumerate}

\vspace{0.5cm}
\noindent\textbf{Nota:} Los criterios técnicos citados a lo largo de la memoria (velocidad objetivo 20~m/s, velocidad máxima de vapor 30~m/s, velocidad en condensados 15--20~m/s, $L_e/D$ de accesorios, pendiente mínima 40~mm/10~m, temperatura superficial máxima 30~°C, profundidad de pozos de goteo 250--700~mm) se fundamentan en las referencias anteriormente listadas, particularmente en el Manual EREN [8], la Guía Spirax Sarco [6], la normativa RITE [1] y la Guía IDAE [5].


